% =====================================================================
%  chuong1.tex  —  CHƯƠNG 1. TỔNG QUAN VỀ ĐỀ TÀI
%  Biên dịch bằng XeLaTeX. Chèn vào main.tex bằng:  % =====================================================================
%  chuong1.tex  —  CHƯƠNG 1. TỔNG QUAN VỀ ĐỀ TÀI
%  Biên dịch bằng XeLaTeX. Chèn vào main.tex bằng:  % =====================================================================
%  chuong1.tex  —  CHƯƠNG 1. TỔNG QUAN VỀ ĐỀ TÀI
%  Biên dịch bằng XeLaTeX. Chèn vào main.tex bằng:  % =====================================================================
%  chuong1.tex  —  CHƯƠNG 1. TỔNG QUAN VỀ ĐỀ TÀI
%  Biên dịch bằng XeLaTeX. Chèn vào main.tex bằng:  \input{chuong1}
%
%  GÓI CẦN CÓ TRONG PREAMBLE:
%     \usepackage{amsmath,amssymb,array}
%     \usepackage[hidelinks]{hyperref}
%     \usepackage[style=apa,backend=biber]{biblatex}
%     \addbibresource{tailieuthamkhao.bib}
%
%  LỆNH TRÍCH DẪN: file này dùng biblatex-apa (\parencite, \textcite).
%     - Nếu khóa luận đang dùng natbib/apacite thì thay:
%           \parencite{key}  ->  \citep{key}
%           \textcite{key}   ->  \citet{key}
%
%  SỐ LIỆU M5 trong mục 1.1.1 và 1.3.2 được tính trực tiếp từ hai tệp
%  calendar.csv và sales_train_evaluation.csv (xem Phụ lục nguồn trích dẫn).
% =====================================================================

\chapter{TỔNG QUAN VỀ ĐỀ TÀI}
\label{chap:tongquan}

\section{Bối cảnh nghiên cứu và tính cấp thiết của đề tài}
\label{sec:boicanh}

\subsection{Thực trạng quản lý tồn kho trong mạng lưới bán lẻ đa điểm}
\label{subsec:thuctrang}

Trong hoạt động bán lẻ hiện đại, hàng hóa không được lưu trữ tập trung tại một
điểm duy nhất mà được phân tán trên một mạng lưới gồm nhiều kho hoặc nhiều cửa
hàng đặt tại các khu vực địa lý khác nhau. Mỗi điểm lưu trữ phục vụ một tập
khách hàng riêng, đối mặt với một chuỗi nhu cầu riêng và phải tự ra quyết định
bổ sung hàng hóa cho hàng trăm đến hàng nghìn mặt hàng khác nhau
\parencite{silver2016}. Bài toán vì vậy không còn là việc xác định một mức đặt
hàng đơn lẻ, mà là việc điều phối đồng thời một số lượng rất lớn các quyết định
có ràng buộc lẫn nhau.

Đặc điểm này thể hiện rõ trong dữ liệu bán lẻ thực tế. Bộ dữ liệu M5, công bố
trong cuộc thi dự báo M5 do \textcite{makridakis2022background} tổ chức, ghi
nhận doanh số bán hàng theo ngày của $3.049$ mặt hàng tại $10$ cửa hàng của
Walmart thuộc ba bang California, Texas và Wisconsin, trải trên $1.941$ ngày từ
ngày 29 tháng 01 năm 2011 đến ngày 22 tháng 05 năm
2016\footnote{Số liệu được tính trực tiếp từ hai tệp \texttt{calendar.csv} và
\texttt{sales\_train\_evaluation.csv} sử dụng trong khóa luận. Việc sử dụng bộ
dữ liệu tuân thủ điều khoản sử dụng dữ liệu của cuộc thi cho mục đích nghiên
cứu học thuật phi thương mại \parencite{m5dataset}.}. Tổ hợp cửa hàng và mặt
hàng tạo thành $30.490$ chuỗi thời gian nhu cầu ở cấp độ chi tiết nhất.

Thống kê trên toàn bộ $30.490$ chuỗi này cho thấy hai đặc trưng có ảnh hưởng
trực tiếp tới thiết kế chính sách đặt hàng. Thứ nhất, nhu cầu ở cấp cửa
hàng--mặt hàng có quy mô rất nhỏ và phân bố lệch mạnh: nhu cầu trung bình là
$1{,}13$ đơn vị mỗi ngày, nhưng trung vị chỉ là $0{,}45$ đơn vị, và $73{,}4\%$
số chuỗi có nhu cầu trung bình dưới một đơn vị mỗi ngày. Thứ hai, nhu cầu mang
tính gián đoạn rõ rệt, với tỷ lệ ngày không phát sinh giao dịch trung bình đạt
$68\%$. Áp dụng tiêu chí phân loại của \textcite{syntetos2005categorization} với
các ngưỡng khoảng cách giữa hai lần phát sinh nhu cầu $p = 1{,}32$ và bình
phương hệ số biến thiên $CV^2 = 0{,}49$, có tới $95{,}2\%$ số chuỗi thuộc hai
nhóm \textit{intermittent} và \textit{lumpy}
(Bảng~\ref{tab:phanloainhucau}).

\begin{table}[htbp]
    \centering
    \caption{Phân loại dạng nhu cầu của $30.490$ chuỗi cửa hàng--mặt hàng trong
    bộ dữ liệu M5 theo tiêu chí của \textcite{syntetos2005categorization}}
    \label{tab:phanloainhucau}
    \begin{tabular}{lcc}
        \hline
        \textbf{Nhóm} & \textbf{Đặc điểm} & \textbf{Tỷ lệ} \\
        \hline
        Smooth       & $p < 1{,}32$; $CV^2 < 0{,}49$   & $3{,}2\%$  \\
        Erratic      & $p < 1{,}32$; $CV^2 \ge 0{,}49$ & $1{,}6\%$  \\
        Intermittent & $p \ge 1{,}32$; $CV^2 < 0{,}49$ & $75{,}7\%$ \\
        Lumpy        & $p \ge 1{,}32$; $CV^2 \ge 0{,}49$ & $19{,}5\%$ \\
        \hline
    \end{tabular}
\end{table}

Đặc điểm gián đoạn này khiến các phương pháp dự báo và các công thức tồn kho
xây dựng trên giả định nhu cầu liên tục và phân phối chuẩn trở nên kém phù hợp
\parencite{croston1972, syntetos2005accuracy}. Quản lý tồn kho trên một mạng
lưới như vậy đòi hỏi những công cụ ra quyết định có khả năng thích ứng, thay vì
các bộ tham số tĩnh được thiết lập một lần và áp dụng cố định.

\subsection{Thách thức trong việc tối ưu hóa chính sách đặt hàng lại}
\label{subsec:thachthuc}

Việc xây dựng chính sách đặt hàng lại trong mạng lưới đa kho -- đa mặt hàng đối
mặt với bốn nhóm thách thức chính.

\textbf{Thứ nhất, tính bất định kép của nhu cầu và thời gian giao hàng.} Nhu
cầu tiêu thụ biến động ngẫu nhiên theo thời gian, đồng thời thời gian từ lúc
phát lệnh đặt hàng đến lúc hàng thực nhập kho cũng không cố định
\parencite{zipkin2000}. Khi thời gian giao hàng là biến ngẫu nhiên, các đơn
hàng phát ra ở những thời điểm khác nhau có thể về kho không đúng thứ tự, khiến
các công thức tính tồn kho an toàn cổ điển chỉ còn giữ vai trò xấp xỉ
\parencite{silver2016}.

\textbf{Thứ hai, bài toán đánh đổi giữa nhiều nhóm chi phí xung đột.} Mục tiêu
vận hành không quy về một đại lượng duy nhất mà là tổng hợp của nhiều thành
phần có chiều tác động ngược nhau, được tóm tắt trong
Bảng~\ref{tab:chiphi}. Giảm tồn kho làm giảm chi phí lưu kho nhưng làm tăng
rủi ro thiếu hàng; đặt hàng theo lô lớn làm giảm chi phí đặt hàng cố định nhưng
làm tăng chi phí lưu kho và nguy cơ vượt sức chứa \parencite{zipkin2000}.

\begin{table}[htbp]
    \centering
    \caption{Các nhóm chi phí trong bài toán tồn kho đang xét}
    \label{tab:chiphi}
    \begin{tabular}{llc}
        \hline
        \textbf{Thành phần chi phí} & \textbf{Ký hiệu} & \textbf{Chiều tác động} \\
        \hline
        Chi phí lưu kho          & $c_{lk}$    & Tăng khi mức tồn kho cao \\
        Chi phí thiếu hàng       & $c_{th}$    & Tăng khi xảy ra hết hàng \\
        Chi phí đặt hàng cố định & $c_{dh}$    & Tăng theo tần suất đặt hàng \\
        Phạt tràn kho            & $p_{tk}$    & Tăng khi vượt sức chứa kho \\
        Phạt vi phạm mức phục vụ & $\phi_{dv}$ & Tăng khi tỷ lệ đáp ứng dưới ngưỡng \\
        \hline
    \end{tabular}
\end{table}

\textbf{Thứ ba, ràng buộc sức chứa tạo ra sự phụ thuộc giữa các mặt hàng.} Các
mặt hàng cùng nằm trong một kho chia sẻ một dung lượng hữu hạn. Quyết định đặt
hàng của một mặt hàng do đó ảnh hưởng trực tiếp tới chi phí mà các mặt hàng
khác trong cùng kho phải gánh chịu. Bài toán vì vậy không thể tách rời thành
các bài toán độc lập, mà mang bản chất của một bài toán ra quyết định hợp tác
\parencite{oliehoek2016}.

\textbf{Thứ tư, quy mô không gian hành động.} Với $N$ kho và $M$ mặt hàng, mỗi
cặp có $K$ mức đặt hàng khả dĩ, không gian hành động chung là tích Descartes
có kích thước $K^{NM}$. Với cấu hình của khóa luận ($N = 10$, $M = 50$,
$K = 6$), con số này là $6^{500}$ --- vượt xa khả năng xử lý của bất kỳ phương
pháp tìm kiếm tập trung nào, và cũng vượt quá khả năng biểu diễn của các thuật
toán học tăng cường dựa trên hàm giá trị vốn phải liệt kê toàn bộ hành động khả
dĩ \parencite{sutton2018}. Trên phương diện lý thuyết độ phức tạp, bài toán
điều khiển phi tập trung dạng này đã được chứng minh thuộc lớp NEXP-đầy đủ
\parencite{bernstein2002}, nên các lời giải chính xác không khả thi ở quy mô
thực tế.

\subsection{Hạn chế của phương pháp truyền thống và sự cần thiết của học tăng
cường đa tác tử}
\label{subsec:khoangtrong}

Các chính sách tồn kho cổ điển đều được xây dựng trên những tập giả định cụ
thể. Mô hình lượng đặt hàng kinh tế của \textcite{harris1913} giả định nhu cầu
xác định và đều, không có tình trạng thiếu hàng. Chính sách $(s, S)$ được
\textcite{scarf1960} chứng minh tối ưu cho hệ đơn kho với nhu cầu độc lập giữa
các kỳ, chi phí đặt hàng cố định và cơ chế giữ đơn hàng chờ. Mô hình người bán
báo của \textcite{arrow1951} chỉ tối ưu cho một chu kỳ đơn lẻ, không xét tới
lượng hàng đang trên đường vận chuyển. Không mô hình nào trong số đó xử lý được
đồng thời nhu cầu ngẫu nhiên, thời gian giao hàng biến động, ràng buộc sức chứa
dùng chung và yêu cầu duy trì mức phục vụ tối thiểu. Riêng với cơ chế mất đơn
hàng khi thiếu hàng --- vốn phù hợp với bối cảnh bán lẻ --- \textcite{zipkin2008}
chỉ ra rằng bài toán không còn nghiệm giải tích dạng đóng và các chính sách cổ
điển mất tính tối ưu đã được chứng minh.

Học tăng cường cung cấp một hướng tiếp cận khác: thay vì giả định trước dạng
phân phối của nhu cầu và suy ra chính sách bằng giải tích, tác tử học chính
sách trực tiếp thông qua tương tác lặp lại với môi trường
\parencite{sutton2018}. Hướng này đã được kiểm chứng trên các bài toán tồn kho
quy mô nhỏ: \textcite{oroojlooyjadid2022} áp dụng mạng Q sâu cho trò chơi phân
phối bia, còn \textcite{gijsbrechts2022} đánh giá học tăng cường sâu trên ba
lớp bài toán tồn kho gồm mất đơn hàng, hai nguồn cung và đa cấp.

Tuy nhiên, việc áp dụng học tăng cường đơn tác tử cho toàn bộ mạng lưới vấp
phải chính vấn đề quy mô đã nêu ở mục trên. Học tăng cường đa tác tử với cơ chế
chia sẻ tham số cho phép vượt qua trở ngại này \parencite{gupta2017}: mỗi cặp
kho--mặt hàng được xem là một tác tử ra quyết định độc lập trên quan sát cục
bộ theo hướng tiếp cận học độc lập của \textcite{tan1993}, trong khi toàn bộ
các tác tử cùng học chung một bộ trọng số. Nhờ đó, không gian hành động mà mỗi
tác tử phải xử lý thu về kích thước $K$, còn số tham số cần học không tăng theo
quy mô hệ thống. Các nghiên cứu gần đây cho thấy phương pháp học độc lập dựa
trên PPO đạt hiệu năng cạnh tranh trên nhiều môi trường hợp tác
\parencite{dewitt2020ippo, yu2022mappo}.

Khảo sát tài liệu cho thấy phần lớn các nghiên cứu áp dụng học tăng cường cho
quản lý tồn kho hiện nay được tiến hành ở quy mô nhỏ, với một vài kho và một
vài mặt hàng, trên dữ liệu nhu cầu sinh theo phân phối lý thuyết. Số công trình
đồng thời mở rộng tới quy mô hàng trăm cặp kho--mặt hàng bằng cơ chế chia sẻ
tham số, đánh giá tường minh ảnh hưởng của cách phân rã hàm phần thưởng ở quy
mô đó, và hiệu chỉnh môi trường mô phỏng từ dữ liệu bán lẻ thực tế còn hạn chế.
Đây chính là khoảng trống mà khóa luận hướng tới.

\section{Mục tiêu và câu hỏi nghiên cứu}
\label{sec:muctieu}

\subsection{Mục tiêu nghiên cứu}
\label{subsec:muctieu}

Khóa luận đặt ra bốn mục tiêu cụ thể:

\begin{itemize}
    \item Xây dựng môi trường mô phỏng tồn kho đa kho -- đa mặt hàng theo chuẩn
    Gymnasium \parencite{towers2024gymnasium}, mô hình hóa đầy đủ nhu cầu ngẫu
    nhiên, thời gian giao hàng biến động, ràng buộc sức chứa và năm nhóm chi phí
    vận hành; tham số nhu cầu được hiệu chỉnh từ bộ dữ liệu M5 Walmart.

    \item Thiết kế hàm phần thưởng phân rã cục bộ kèm thành phần phạt dựa trên
    tỷ lệ đáp ứng nhu cầu tính trên cửa sổ trượt, bảo đảm tín hiệu học phản ánh
    đúng đóng góp của từng tác tử \parencite{wolpert2001}.

    \item Cài đặt và huấn luyện thuật toán Independent Multi-Agent PPO
    \parencite{schulman2017ppo, dewitt2020ippo} với cơ chế chia sẻ tham số cho
    $500$ cặp kho--mặt hàng, sử dụng ước lượng lợi thế tổng quát
    \parencite{schulman2016gae} và chuẩn hóa lợi thế theo từng cặp.

    \item Đánh giá hiệu năng của mô hình đề xuất trong so sánh với ba chính
    sách tồn kho cổ điển, theo giao thức thực nghiệm đa hạt giống có kiểm định
    thống kê \parencite{henderson2018, agarwal2021}.
\end{itemize}

\subsection{Câu hỏi nghiên cứu}
\label{subsec:cauhoi}

Khóa luận hướng tới trả lời ba câu hỏi:

\begin{enumerate}
    \item \textbf{Câu hỏi 1.} Tác tử IPPO có học được chính sách đặt hàng lại
    cho tổng chi phí vận hành thấp hơn các chính sách tối ưu hóa cổ điển
    (EOQ, $(s, S)$, Newsvendor) trong môi trường đa kho -- đa mặt hàng với nhu
    cầu ngẫu nhiên và thời gian giao hàng biến động hay không, và mức cải thiện
    đó có ý nghĩa thống kê hay không?

    \item \textbf{Câu hỏi 2.} Cơ chế chia sẻ tham số giữa $500$ tác tử có đủ để
    học một chính sách tổng quát hóa được cho các mặt hàng có quy mô nhu cầu
    khác biệt đáng kể hay không?

    \item \textbf{Câu hỏi 3.} Hàm phần thưởng phân rã cục bộ kèm tỷ lệ đáp ứng
    theo cửa sổ trượt có cải thiện quá trình học so với hàm phần thưởng toàn
    cục dùng chung hay không?
\end{enumerate}

\section{Đối tượng và phạm vi nghiên cứu}
\label{sec:phamvi}

\subsection{Đối tượng nghiên cứu}
\label{subsec:doituong}

Đối tượng nghiên cứu của khóa luận là bài toán tối ưu hóa chính sách đặt hàng
lại trong hệ thống tồn kho đa kho -- đa mặt hàng, và thuật toán học tăng cường
đa tác tử Independent Multi-Agent PPO áp dụng cho bài toán này.

\subsection{Phạm vi trong đề tài}
\label{subsec:phamvitrong}

\begin{itemize}
    \item \textbf{Cấu hình hệ thống:} $N = 10$ kho, $M = 50$ mặt hàng, tương
    ứng $500$ cặp ra quyết định; độ dài mỗi tập huấn luyện là $365$ ngày; thời
    gian giao hàng ngẫu nhiên trong khoảng từ một đến ba ngày.

    \item \textbf{Cấu trúc mạng lưới:} các kho hoạt động song song và độc lập
    về nguồn cung, mỗi kho đặt hàng trực tiếp từ nhà cung cấp bên ngoài. Sự
    ràng buộc giữa các mặt hàng phát sinh từ dung lượng kho dùng chung.

    \item \textbf{Nguồn dữ liệu:} bộ dữ liệu bán lẻ M5 Walmart
    \parencite{makridakis2022background}, cụ thể là hai tệp
    \texttt{sales\_train\_evaluation.csv} chứa $30.490$ chuỗi doanh số theo
    ngày trên $1.941$ ngày, và \texttt{calendar.csv} chứa thông tin lịch, ngày
    lễ và các sự kiện trợ cấp thực phẩm SNAP. Khóa luận đồng thời hỗ trợ sinh
    dữ liệu giả lập theo phân phối Poisson phục vụ thử nghiệm đối chứng.

    \item \textbf{Phương pháp đối chứng:} ba chính sách tồn kho cổ điển EOQ
    \parencite{harris1913}, $(s, S)$ \parencite{scarf1960} và Newsvendor
    \parencite{arrow1951}, với tham số được tối ưu bằng tìm kiếm lưới trên
    chính môi trường mô phỏng nhằm bảo đảm tính công bằng của so sánh.

    \item \textbf{Chỉ số đánh giá:} tổng chi phí vận hành, tỷ lệ đáp ứng nhu
    cầu, và cơ cấu các thành phần chi phí lưu kho, thiếu hàng, đặt hàng.
\end{itemize}

\subsection{Giới hạn phạm vi}
\label{subsec:phamvingoai}

Để bảo đảm chiều sâu nghiên cứu trong thời gian thực hiện, những nội dung sau
được xác định nằm ngoài phạm vi khóa luận:

\begin{itemize}
    \item Cấu trúc tồn kho đa cấp (multi-echelon) với quan hệ cung ứng giữa các
    kho theo mô hình của \textcite{clark1960}, cùng các hiện tượng đi kèm như
    hiệu ứng bullwhip \parencite{lee1997}.
    \item Bài toán điều chuyển hàng liên kho và tối ưu hóa tuyến vận chuyển.
    \item Dự báo nhu cầu bằng mô hình học sâu chuỗi thời gian; khóa luận sử
    dụng trực tiếp thông tin nhu cầu lịch sử trong vector quan sát.
    \item Cơ chế truyền thông giữa các tác tử; các tác tử ra quyết định hoàn
    toàn dựa trên quan sát cục bộ.
    \item Triển khai hệ thống thời gian thực hoặc tích hợp với nền tảng quản trị
    doanh nghiệp; sản phẩm của khóa luận là kết quả mô phỏng và phân tích định
    lượng.
\end{itemize}

\section{Phương pháp nghiên cứu}
\label{sec:phuongphap}

Khóa luận kết hợp ba nhóm phương pháp. Về \textbf{lý thuyết}, bài toán được mô
hình hóa dưới dạng quá trình quyết định Markov phi tập trung có quan sát bộ
phận \parencite{bernstein2002, oliehoek2016}, trên nền tảng lý thuyết kiểm soát
tồn kho ngẫu nhiên \parencite{zipkin2000} và học tăng cường
\parencite{puterman1994, sutton2018}. Về \textbf{thực nghiệm}, nhóm xây dựng
môi trường mô phỏng có kiểm thử đơn vị theo chuẩn giao diện Gymnasium
\parencite{brockman2016gym, towers2024gymnasium}, hiệu chỉnh tham số từ dữ liệu
M5, huấn luyện tác tử và thu thập số liệu vận hành. Về \textbf{phân tích}, mọi
so sánh giữa các phương pháp đều được thực hiện trên tối thiểu năm hạt giống
ngẫu nhiên, báo cáo giá trị trung bình kèm độ lệch chuẩn, và được xác nhận bằng
kiểm định thống kê phù hợp trước khi kết luận --- theo khuyến nghị về tính tái
lập trong đánh giá học tăng cường sâu của \textcite{henderson2018} và
\textcite{agarwal2021}.

\section{Ý nghĩa khoa học và thực tiễn của đề tài}
\label{sec:ynghia}

\textbf{Về mặt khoa học}, khóa luận đóng góp một khung mô hình hóa bài toán tồn
kho đa kho -- đa mặt hàng dưới dạng hệ đa tác tử hợp tác quy mô lớn, cùng với
đánh giá thực nghiệm có kiểm định thống kê về ảnh hưởng của cách phân rã hàm
phần thưởng và cơ chế chia sẻ tham số --- hai yếu tố thiết kế thường được sử
dụng nhưng ít được khảo sát tách bạch trong các nghiên cứu trước.

\textbf{Về mặt thực tiễn}, kết quả của khóa luận cung cấp bằng chứng định lượng
về mức độ khả thi của việc thay thế các chính sách đặt hàng theo tham số tĩnh
bằng chính sách học được từ dữ liệu, trên một môi trường mô phỏng được hiệu
chỉnh từ dữ liệu bán lẻ thực tế. Cấu trúc mã nguồn và giao thức đánh giá được
tổ chức theo hướng có thể tái lập, tạo điều kiện cho các nghiên cứu mở rộng
tiếp theo.

\section{Bố cục của khóa luận}
\label{sec:bocuc}

Khóa luận được tổ chức thành năm chương:

\begin{itemize}
    \item \textbf{Chương 1 --- Tổng quan về đề tài:} trình bày bối cảnh, tính
    cấp thiết, mục tiêu, câu hỏi nghiên cứu, phạm vi và phương pháp nghiên cứu.

    \item \textbf{Chương 2 --- Cơ sở lý thuyết và các công trình liên quan:}
    hệ thống hóa lý thuyết kiểm soát tồn kho ngẫu nhiên, các chính sách cổ
    điển, nền tảng học tăng cường và học tăng cường đa tác tử, cùng lược khảo
    các nghiên cứu tiền nhiệm.

    \item \textbf{Chương 3 --- Phân tích và thiết kế mô hình:} chi tiết hóa mô
    hình toán học của môi trường, thiết kế không gian trạng thái, không gian
    hành động và hàm phần thưởng, kiến trúc mạng Actor--Critic chia sẻ tham số
    và quy trình huấn luyện.

    \item \textbf{Chương 4 --- Thực nghiệm và đánh giá kết quả:} trình bày cấu
    hình thực nghiệm, kết quả so sánh với các chính sách cổ điển, các khảo sát
    loại trừ và phân tích độ nhạy.

    \item \textbf{Chương 5 --- Kết luận và hướng phát triển:} tổng kết kết quả
    đạt được, chỉ ra hạn chế và đề xuất hướng nghiên cứu tiếp theo.
\end{itemize}

%
%  GÓI CẦN CÓ TRONG PREAMBLE:
%     \usepackage{amsmath,amssymb,array}
%     \usepackage[hidelinks]{hyperref}
%     \usepackage[style=apa,backend=biber]{biblatex}
%     \addbibresource{tailieuthamkhao.bib}
%
%  LỆNH TRÍCH DẪN: file này dùng biblatex-apa (\parencite, \textcite).
%     - Nếu khóa luận đang dùng natbib/apacite thì thay:
%           \parencite{key}  ->  \citep{key}
%           \textcite{key}   ->  \citet{key}
%
%  SỐ LIỆU M5 trong mục 1.1.1 và 1.3.2 được tính trực tiếp từ hai tệp
%  calendar.csv và sales_train_evaluation.csv (xem Phụ lục nguồn trích dẫn).
% =====================================================================

\chapter{TỔNG QUAN VỀ ĐỀ TÀI}
\label{chap:tongquan}

\section{Bối cảnh nghiên cứu và tính cấp thiết của đề tài}
\label{sec:boicanh}

\subsection{Thực trạng quản lý tồn kho trong mạng lưới bán lẻ đa điểm}
\label{subsec:thuctrang}

Trong hoạt động bán lẻ hiện đại, hàng hóa không được lưu trữ tập trung tại một
điểm duy nhất mà được phân tán trên một mạng lưới gồm nhiều kho hoặc nhiều cửa
hàng đặt tại các khu vực địa lý khác nhau. Mỗi điểm lưu trữ phục vụ một tập
khách hàng riêng, đối mặt với một chuỗi nhu cầu riêng và phải tự ra quyết định
bổ sung hàng hóa cho hàng trăm đến hàng nghìn mặt hàng khác nhau
\parencite{silver2016}. Bài toán vì vậy không còn là việc xác định một mức đặt
hàng đơn lẻ, mà là việc điều phối đồng thời một số lượng rất lớn các quyết định
có ràng buộc lẫn nhau.

Đặc điểm này thể hiện rõ trong dữ liệu bán lẻ thực tế. Bộ dữ liệu M5, công bố
trong cuộc thi dự báo M5 do \textcite{makridakis2022background} tổ chức, ghi
nhận doanh số bán hàng theo ngày của $3.049$ mặt hàng tại $10$ cửa hàng của
Walmart thuộc ba bang California, Texas và Wisconsin, trải trên $1.941$ ngày từ
ngày 29 tháng 01 năm 2011 đến ngày 22 tháng 05 năm
2016\footnote{Số liệu được tính trực tiếp từ hai tệp \texttt{calendar.csv} và
\texttt{sales\_train\_evaluation.csv} sử dụng trong khóa luận. Việc sử dụng bộ
dữ liệu tuân thủ điều khoản sử dụng dữ liệu của cuộc thi cho mục đích nghiên
cứu học thuật phi thương mại \parencite{m5dataset}.}. Tổ hợp cửa hàng và mặt
hàng tạo thành $30.490$ chuỗi thời gian nhu cầu ở cấp độ chi tiết nhất.

Thống kê trên toàn bộ $30.490$ chuỗi này cho thấy hai đặc trưng có ảnh hưởng
trực tiếp tới thiết kế chính sách đặt hàng. Thứ nhất, nhu cầu ở cấp cửa
hàng--mặt hàng có quy mô rất nhỏ và phân bố lệch mạnh: nhu cầu trung bình là
$1{,}13$ đơn vị mỗi ngày, nhưng trung vị chỉ là $0{,}45$ đơn vị, và $73{,}4\%$
số chuỗi có nhu cầu trung bình dưới một đơn vị mỗi ngày. Thứ hai, nhu cầu mang
tính gián đoạn rõ rệt, với tỷ lệ ngày không phát sinh giao dịch trung bình đạt
$68\%$. Áp dụng tiêu chí phân loại của \textcite{syntetos2005categorization} với
các ngưỡng khoảng cách giữa hai lần phát sinh nhu cầu $p = 1{,}32$ và bình
phương hệ số biến thiên $CV^2 = 0{,}49$, có tới $95{,}2\%$ số chuỗi thuộc hai
nhóm \textit{intermittent} và \textit{lumpy}
(Bảng~\ref{tab:phanloainhucau}).

\begin{table}[htbp]
    \centering
    \caption{Phân loại dạng nhu cầu của $30.490$ chuỗi cửa hàng--mặt hàng trong
    bộ dữ liệu M5 theo tiêu chí của \textcite{syntetos2005categorization}}
    \label{tab:phanloainhucau}
    \begin{tabular}{lcc}
        \hline
        \textbf{Nhóm} & \textbf{Đặc điểm} & \textbf{Tỷ lệ} \\
        \hline
        Smooth       & $p < 1{,}32$; $CV^2 < 0{,}49$   & $3{,}2\%$  \\
        Erratic      & $p < 1{,}32$; $CV^2 \ge 0{,}49$ & $1{,}6\%$  \\
        Intermittent & $p \ge 1{,}32$; $CV^2 < 0{,}49$ & $75{,}7\%$ \\
        Lumpy        & $p \ge 1{,}32$; $CV^2 \ge 0{,}49$ & $19{,}5\%$ \\
        \hline
    \end{tabular}
\end{table}

Đặc điểm gián đoạn này khiến các phương pháp dự báo và các công thức tồn kho
xây dựng trên giả định nhu cầu liên tục và phân phối chuẩn trở nên kém phù hợp
\parencite{croston1972, syntetos2005accuracy}. Quản lý tồn kho trên một mạng
lưới như vậy đòi hỏi những công cụ ra quyết định có khả năng thích ứng, thay vì
các bộ tham số tĩnh được thiết lập một lần và áp dụng cố định.

\subsection{Thách thức trong việc tối ưu hóa chính sách đặt hàng lại}
\label{subsec:thachthuc}

Việc xây dựng chính sách đặt hàng lại trong mạng lưới đa kho -- đa mặt hàng đối
mặt với bốn nhóm thách thức chính.

\textbf{Thứ nhất, tính bất định kép của nhu cầu và thời gian giao hàng.} Nhu
cầu tiêu thụ biến động ngẫu nhiên theo thời gian, đồng thời thời gian từ lúc
phát lệnh đặt hàng đến lúc hàng thực nhập kho cũng không cố định
\parencite{zipkin2000}. Khi thời gian giao hàng là biến ngẫu nhiên, các đơn
hàng phát ra ở những thời điểm khác nhau có thể về kho không đúng thứ tự, khiến
các công thức tính tồn kho an toàn cổ điển chỉ còn giữ vai trò xấp xỉ
\parencite{silver2016}.

\textbf{Thứ hai, bài toán đánh đổi giữa nhiều nhóm chi phí xung đột.} Mục tiêu
vận hành không quy về một đại lượng duy nhất mà là tổng hợp của nhiều thành
phần có chiều tác động ngược nhau, được tóm tắt trong
Bảng~\ref{tab:chiphi}. Giảm tồn kho làm giảm chi phí lưu kho nhưng làm tăng
rủi ro thiếu hàng; đặt hàng theo lô lớn làm giảm chi phí đặt hàng cố định nhưng
làm tăng chi phí lưu kho và nguy cơ vượt sức chứa \parencite{zipkin2000}.

\begin{table}[htbp]
    \centering
    \caption{Các nhóm chi phí trong bài toán tồn kho đang xét}
    \label{tab:chiphi}
    \begin{tabular}{llc}
        \hline
        \textbf{Thành phần chi phí} & \textbf{Ký hiệu} & \textbf{Chiều tác động} \\
        \hline
        Chi phí lưu kho          & $c_{lk}$    & Tăng khi mức tồn kho cao \\
        Chi phí thiếu hàng       & $c_{th}$    & Tăng khi xảy ra hết hàng \\
        Chi phí đặt hàng cố định & $c_{dh}$    & Tăng theo tần suất đặt hàng \\
        Phạt tràn kho            & $p_{tk}$    & Tăng khi vượt sức chứa kho \\
        Phạt vi phạm mức phục vụ & $\phi_{dv}$ & Tăng khi tỷ lệ đáp ứng dưới ngưỡng \\
        \hline
    \end{tabular}
\end{table}

\textbf{Thứ ba, ràng buộc sức chứa tạo ra sự phụ thuộc giữa các mặt hàng.} Các
mặt hàng cùng nằm trong một kho chia sẻ một dung lượng hữu hạn. Quyết định đặt
hàng của một mặt hàng do đó ảnh hưởng trực tiếp tới chi phí mà các mặt hàng
khác trong cùng kho phải gánh chịu. Bài toán vì vậy không thể tách rời thành
các bài toán độc lập, mà mang bản chất của một bài toán ra quyết định hợp tác
\parencite{oliehoek2016}.

\textbf{Thứ tư, quy mô không gian hành động.} Với $N$ kho và $M$ mặt hàng, mỗi
cặp có $K$ mức đặt hàng khả dĩ, không gian hành động chung là tích Descartes
có kích thước $K^{NM}$. Với cấu hình của khóa luận ($N = 10$, $M = 50$,
$K = 6$), con số này là $6^{500}$ --- vượt xa khả năng xử lý của bất kỳ phương
pháp tìm kiếm tập trung nào, và cũng vượt quá khả năng biểu diễn của các thuật
toán học tăng cường dựa trên hàm giá trị vốn phải liệt kê toàn bộ hành động khả
dĩ \parencite{sutton2018}. Trên phương diện lý thuyết độ phức tạp, bài toán
điều khiển phi tập trung dạng này đã được chứng minh thuộc lớp NEXP-đầy đủ
\parencite{bernstein2002}, nên các lời giải chính xác không khả thi ở quy mô
thực tế.

\subsection{Hạn chế của phương pháp truyền thống và sự cần thiết của học tăng
cường đa tác tử}
\label{subsec:khoangtrong}

Các chính sách tồn kho cổ điển đều được xây dựng trên những tập giả định cụ
thể. Mô hình lượng đặt hàng kinh tế của \textcite{harris1913} giả định nhu cầu
xác định và đều, không có tình trạng thiếu hàng. Chính sách $(s, S)$ được
\textcite{scarf1960} chứng minh tối ưu cho hệ đơn kho với nhu cầu độc lập giữa
các kỳ, chi phí đặt hàng cố định và cơ chế giữ đơn hàng chờ. Mô hình người bán
báo của \textcite{arrow1951} chỉ tối ưu cho một chu kỳ đơn lẻ, không xét tới
lượng hàng đang trên đường vận chuyển. Không mô hình nào trong số đó xử lý được
đồng thời nhu cầu ngẫu nhiên, thời gian giao hàng biến động, ràng buộc sức chứa
dùng chung và yêu cầu duy trì mức phục vụ tối thiểu. Riêng với cơ chế mất đơn
hàng khi thiếu hàng --- vốn phù hợp với bối cảnh bán lẻ --- \textcite{zipkin2008}
chỉ ra rằng bài toán không còn nghiệm giải tích dạng đóng và các chính sách cổ
điển mất tính tối ưu đã được chứng minh.

Học tăng cường cung cấp một hướng tiếp cận khác: thay vì giả định trước dạng
phân phối của nhu cầu và suy ra chính sách bằng giải tích, tác tử học chính
sách trực tiếp thông qua tương tác lặp lại với môi trường
\parencite{sutton2018}. Hướng này đã được kiểm chứng trên các bài toán tồn kho
quy mô nhỏ: \textcite{oroojlooyjadid2022} áp dụng mạng Q sâu cho trò chơi phân
phối bia, còn \textcite{gijsbrechts2022} đánh giá học tăng cường sâu trên ba
lớp bài toán tồn kho gồm mất đơn hàng, hai nguồn cung và đa cấp.

Tuy nhiên, việc áp dụng học tăng cường đơn tác tử cho toàn bộ mạng lưới vấp
phải chính vấn đề quy mô đã nêu ở mục trên. Học tăng cường đa tác tử với cơ chế
chia sẻ tham số cho phép vượt qua trở ngại này \parencite{gupta2017}: mỗi cặp
kho--mặt hàng được xem là một tác tử ra quyết định độc lập trên quan sát cục
bộ theo hướng tiếp cận học độc lập của \textcite{tan1993}, trong khi toàn bộ
các tác tử cùng học chung một bộ trọng số. Nhờ đó, không gian hành động mà mỗi
tác tử phải xử lý thu về kích thước $K$, còn số tham số cần học không tăng theo
quy mô hệ thống. Các nghiên cứu gần đây cho thấy phương pháp học độc lập dựa
trên PPO đạt hiệu năng cạnh tranh trên nhiều môi trường hợp tác
\parencite{dewitt2020ippo, yu2022mappo}.

Khảo sát tài liệu cho thấy phần lớn các nghiên cứu áp dụng học tăng cường cho
quản lý tồn kho hiện nay được tiến hành ở quy mô nhỏ, với một vài kho và một
vài mặt hàng, trên dữ liệu nhu cầu sinh theo phân phối lý thuyết. Số công trình
đồng thời mở rộng tới quy mô hàng trăm cặp kho--mặt hàng bằng cơ chế chia sẻ
tham số, đánh giá tường minh ảnh hưởng của cách phân rã hàm phần thưởng ở quy
mô đó, và hiệu chỉnh môi trường mô phỏng từ dữ liệu bán lẻ thực tế còn hạn chế.
Đây chính là khoảng trống mà khóa luận hướng tới.

\section{Mục tiêu và câu hỏi nghiên cứu}
\label{sec:muctieu}

\subsection{Mục tiêu nghiên cứu}
\label{subsec:muctieu}

Khóa luận đặt ra bốn mục tiêu cụ thể:

\begin{itemize}
    \item Xây dựng môi trường mô phỏng tồn kho đa kho -- đa mặt hàng theo chuẩn
    Gymnasium \parencite{towers2024gymnasium}, mô hình hóa đầy đủ nhu cầu ngẫu
    nhiên, thời gian giao hàng biến động, ràng buộc sức chứa và năm nhóm chi phí
    vận hành; tham số nhu cầu được hiệu chỉnh từ bộ dữ liệu M5 Walmart.

    \item Thiết kế hàm phần thưởng phân rã cục bộ kèm thành phần phạt dựa trên
    tỷ lệ đáp ứng nhu cầu tính trên cửa sổ trượt, bảo đảm tín hiệu học phản ánh
    đúng đóng góp của từng tác tử \parencite{wolpert2001}.

    \item Cài đặt và huấn luyện thuật toán Independent Multi-Agent PPO
    \parencite{schulman2017ppo, dewitt2020ippo} với cơ chế chia sẻ tham số cho
    $500$ cặp kho--mặt hàng, sử dụng ước lượng lợi thế tổng quát
    \parencite{schulman2016gae} và chuẩn hóa lợi thế theo từng cặp.

    \item Đánh giá hiệu năng của mô hình đề xuất trong so sánh với ba chính
    sách tồn kho cổ điển, theo giao thức thực nghiệm đa hạt giống có kiểm định
    thống kê \parencite{henderson2018, agarwal2021}.
\end{itemize}

\subsection{Câu hỏi nghiên cứu}
\label{subsec:cauhoi}

Khóa luận hướng tới trả lời ba câu hỏi:

\begin{enumerate}
    \item \textbf{Câu hỏi 1.} Tác tử IPPO có học được chính sách đặt hàng lại
    cho tổng chi phí vận hành thấp hơn các chính sách tối ưu hóa cổ điển
    (EOQ, $(s, S)$, Newsvendor) trong môi trường đa kho -- đa mặt hàng với nhu
    cầu ngẫu nhiên và thời gian giao hàng biến động hay không, và mức cải thiện
    đó có ý nghĩa thống kê hay không?

    \item \textbf{Câu hỏi 2.} Cơ chế chia sẻ tham số giữa $500$ tác tử có đủ để
    học một chính sách tổng quát hóa được cho các mặt hàng có quy mô nhu cầu
    khác biệt đáng kể hay không?

    \item \textbf{Câu hỏi 3.} Hàm phần thưởng phân rã cục bộ kèm tỷ lệ đáp ứng
    theo cửa sổ trượt có cải thiện quá trình học so với hàm phần thưởng toàn
    cục dùng chung hay không?
\end{enumerate}

\section{Đối tượng và phạm vi nghiên cứu}
\label{sec:phamvi}

\subsection{Đối tượng nghiên cứu}
\label{subsec:doituong}

Đối tượng nghiên cứu của khóa luận là bài toán tối ưu hóa chính sách đặt hàng
lại trong hệ thống tồn kho đa kho -- đa mặt hàng, và thuật toán học tăng cường
đa tác tử Independent Multi-Agent PPO áp dụng cho bài toán này.

\subsection{Phạm vi trong đề tài}
\label{subsec:phamvitrong}

\begin{itemize}
    \item \textbf{Cấu hình hệ thống:} $N = 10$ kho, $M = 50$ mặt hàng, tương
    ứng $500$ cặp ra quyết định; độ dài mỗi tập huấn luyện là $365$ ngày; thời
    gian giao hàng ngẫu nhiên trong khoảng từ một đến ba ngày.

    \item \textbf{Cấu trúc mạng lưới:} các kho hoạt động song song và độc lập
    về nguồn cung, mỗi kho đặt hàng trực tiếp từ nhà cung cấp bên ngoài. Sự
    ràng buộc giữa các mặt hàng phát sinh từ dung lượng kho dùng chung.

    \item \textbf{Nguồn dữ liệu:} bộ dữ liệu bán lẻ M5 Walmart
    \parencite{makridakis2022background}, cụ thể là hai tệp
    \texttt{sales\_train\_evaluation.csv} chứa $30.490$ chuỗi doanh số theo
    ngày trên $1.941$ ngày, và \texttt{calendar.csv} chứa thông tin lịch, ngày
    lễ và các sự kiện trợ cấp thực phẩm SNAP. Khóa luận đồng thời hỗ trợ sinh
    dữ liệu giả lập theo phân phối Poisson phục vụ thử nghiệm đối chứng.

    \item \textbf{Phương pháp đối chứng:} ba chính sách tồn kho cổ điển EOQ
    \parencite{harris1913}, $(s, S)$ \parencite{scarf1960} và Newsvendor
    \parencite{arrow1951}, với tham số được tối ưu bằng tìm kiếm lưới trên
    chính môi trường mô phỏng nhằm bảo đảm tính công bằng của so sánh.

    \item \textbf{Chỉ số đánh giá:} tổng chi phí vận hành, tỷ lệ đáp ứng nhu
    cầu, và cơ cấu các thành phần chi phí lưu kho, thiếu hàng, đặt hàng.
\end{itemize}

\subsection{Giới hạn phạm vi}
\label{subsec:phamvingoai}

Để bảo đảm chiều sâu nghiên cứu trong thời gian thực hiện, những nội dung sau
được xác định nằm ngoài phạm vi khóa luận:

\begin{itemize}
    \item Cấu trúc tồn kho đa cấp (multi-echelon) với quan hệ cung ứng giữa các
    kho theo mô hình của \textcite{clark1960}, cùng các hiện tượng đi kèm như
    hiệu ứng bullwhip \parencite{lee1997}.
    \item Bài toán điều chuyển hàng liên kho và tối ưu hóa tuyến vận chuyển.
    \item Dự báo nhu cầu bằng mô hình học sâu chuỗi thời gian; khóa luận sử
    dụng trực tiếp thông tin nhu cầu lịch sử trong vector quan sát.
    \item Cơ chế truyền thông giữa các tác tử; các tác tử ra quyết định hoàn
    toàn dựa trên quan sát cục bộ.
    \item Triển khai hệ thống thời gian thực hoặc tích hợp với nền tảng quản trị
    doanh nghiệp; sản phẩm của khóa luận là kết quả mô phỏng và phân tích định
    lượng.
\end{itemize}

\section{Phương pháp nghiên cứu}
\label{sec:phuongphap}

Khóa luận kết hợp ba nhóm phương pháp. Về \textbf{lý thuyết}, bài toán được mô
hình hóa dưới dạng quá trình quyết định Markov phi tập trung có quan sát bộ
phận \parencite{bernstein2002, oliehoek2016}, trên nền tảng lý thuyết kiểm soát
tồn kho ngẫu nhiên \parencite{zipkin2000} và học tăng cường
\parencite{puterman1994, sutton2018}. Về \textbf{thực nghiệm}, nhóm xây dựng
môi trường mô phỏng có kiểm thử đơn vị theo chuẩn giao diện Gymnasium
\parencite{brockman2016gym, towers2024gymnasium}, hiệu chỉnh tham số từ dữ liệu
M5, huấn luyện tác tử và thu thập số liệu vận hành. Về \textbf{phân tích}, mọi
so sánh giữa các phương pháp đều được thực hiện trên tối thiểu năm hạt giống
ngẫu nhiên, báo cáo giá trị trung bình kèm độ lệch chuẩn, và được xác nhận bằng
kiểm định thống kê phù hợp trước khi kết luận --- theo khuyến nghị về tính tái
lập trong đánh giá học tăng cường sâu của \textcite{henderson2018} và
\textcite{agarwal2021}.

\section{Ý nghĩa khoa học và thực tiễn của đề tài}
\label{sec:ynghia}

\textbf{Về mặt khoa học}, khóa luận đóng góp một khung mô hình hóa bài toán tồn
kho đa kho -- đa mặt hàng dưới dạng hệ đa tác tử hợp tác quy mô lớn, cùng với
đánh giá thực nghiệm có kiểm định thống kê về ảnh hưởng của cách phân rã hàm
phần thưởng và cơ chế chia sẻ tham số --- hai yếu tố thiết kế thường được sử
dụng nhưng ít được khảo sát tách bạch trong các nghiên cứu trước.

\textbf{Về mặt thực tiễn}, kết quả của khóa luận cung cấp bằng chứng định lượng
về mức độ khả thi của việc thay thế các chính sách đặt hàng theo tham số tĩnh
bằng chính sách học được từ dữ liệu, trên một môi trường mô phỏng được hiệu
chỉnh từ dữ liệu bán lẻ thực tế. Cấu trúc mã nguồn và giao thức đánh giá được
tổ chức theo hướng có thể tái lập, tạo điều kiện cho các nghiên cứu mở rộng
tiếp theo.

\section{Bố cục của khóa luận}
\label{sec:bocuc}

Khóa luận được tổ chức thành năm chương:

\begin{itemize}
    \item \textbf{Chương 1 --- Tổng quan về đề tài:} trình bày bối cảnh, tính
    cấp thiết, mục tiêu, câu hỏi nghiên cứu, phạm vi và phương pháp nghiên cứu.

    \item \textbf{Chương 2 --- Cơ sở lý thuyết và các công trình liên quan:}
    hệ thống hóa lý thuyết kiểm soát tồn kho ngẫu nhiên, các chính sách cổ
    điển, nền tảng học tăng cường và học tăng cường đa tác tử, cùng lược khảo
    các nghiên cứu tiền nhiệm.

    \item \textbf{Chương 3 --- Phân tích và thiết kế mô hình:} chi tiết hóa mô
    hình toán học của môi trường, thiết kế không gian trạng thái, không gian
    hành động và hàm phần thưởng, kiến trúc mạng Actor--Critic chia sẻ tham số
    và quy trình huấn luyện.

    \item \textbf{Chương 4 --- Thực nghiệm và đánh giá kết quả:} trình bày cấu
    hình thực nghiệm, kết quả so sánh với các chính sách cổ điển, các khảo sát
    loại trừ và phân tích độ nhạy.

    \item \textbf{Chương 5 --- Kết luận và hướng phát triển:} tổng kết kết quả
    đạt được, chỉ ra hạn chế và đề xuất hướng nghiên cứu tiếp theo.
\end{itemize}

%
%  GÓI CẦN CÓ TRONG PREAMBLE:
%     \usepackage{amsmath,amssymb,array}
%     \usepackage[hidelinks]{hyperref}
%     \usepackage[style=apa,backend=biber]{biblatex}
%     \addbibresource{tailieuthamkhao.bib}
%
%  LỆNH TRÍCH DẪN: file này dùng biblatex-apa (\parencite, \textcite).
%     - Nếu khóa luận đang dùng natbib/apacite thì thay:
%           \parencite{key}  ->  \citep{key}
%           \textcite{key}   ->  \citet{key}
%
%  SỐ LIỆU M5 trong mục 1.1.1 và 1.3.2 được tính trực tiếp từ hai tệp
%  calendar.csv và sales_train_evaluation.csv (xem Phụ lục nguồn trích dẫn).
% =====================================================================

\chapter{TỔNG QUAN VỀ ĐỀ TÀI}
\label{chap:tongquan}

\section{Bối cảnh nghiên cứu và tính cấp thiết của đề tài}
\label{sec:boicanh}

\subsection{Thực trạng quản lý tồn kho trong mạng lưới bán lẻ đa điểm}
\label{subsec:thuctrang}

Trong hoạt động bán lẻ hiện đại, hàng hóa không được lưu trữ tập trung tại một
điểm duy nhất mà được phân tán trên một mạng lưới gồm nhiều kho hoặc nhiều cửa
hàng đặt tại các khu vực địa lý khác nhau. Mỗi điểm lưu trữ phục vụ một tập
khách hàng riêng, đối mặt với một chuỗi nhu cầu riêng và phải tự ra quyết định
bổ sung hàng hóa cho hàng trăm đến hàng nghìn mặt hàng khác nhau
\parencite{silver2016}. Bài toán vì vậy không còn là việc xác định một mức đặt
hàng đơn lẻ, mà là việc điều phối đồng thời một số lượng rất lớn các quyết định
có ràng buộc lẫn nhau.

Đặc điểm này thể hiện rõ trong dữ liệu bán lẻ thực tế. Bộ dữ liệu M5, công bố
trong cuộc thi dự báo M5 do \textcite{makridakis2022background} tổ chức, ghi
nhận doanh số bán hàng theo ngày của $3.049$ mặt hàng tại $10$ cửa hàng của
Walmart thuộc ba bang California, Texas và Wisconsin, trải trên $1.941$ ngày từ
ngày 29 tháng 01 năm 2011 đến ngày 22 tháng 05 năm
2016\footnote{Số liệu được tính trực tiếp từ hai tệp \texttt{calendar.csv} và
\texttt{sales\_train\_evaluation.csv} sử dụng trong khóa luận. Việc sử dụng bộ
dữ liệu tuân thủ điều khoản sử dụng dữ liệu của cuộc thi cho mục đích nghiên
cứu học thuật phi thương mại \parencite{m5dataset}.}. Tổ hợp cửa hàng và mặt
hàng tạo thành $30.490$ chuỗi thời gian nhu cầu ở cấp độ chi tiết nhất.

Thống kê trên toàn bộ $30.490$ chuỗi này cho thấy hai đặc trưng có ảnh hưởng
trực tiếp tới thiết kế chính sách đặt hàng. Thứ nhất, nhu cầu ở cấp cửa
hàng--mặt hàng có quy mô rất nhỏ và phân bố lệch mạnh: nhu cầu trung bình là
$1{,}13$ đơn vị mỗi ngày, nhưng trung vị chỉ là $0{,}45$ đơn vị, và $73{,}4\%$
số chuỗi có nhu cầu trung bình dưới một đơn vị mỗi ngày. Thứ hai, nhu cầu mang
tính gián đoạn rõ rệt, với tỷ lệ ngày không phát sinh giao dịch trung bình đạt
$68\%$. Áp dụng tiêu chí phân loại của \textcite{syntetos2005categorization} với
các ngưỡng khoảng cách giữa hai lần phát sinh nhu cầu $p = 1{,}32$ và bình
phương hệ số biến thiên $CV^2 = 0{,}49$, có tới $95{,}2\%$ số chuỗi thuộc hai
nhóm \textit{intermittent} và \textit{lumpy}
(Bảng~\ref{tab:phanloainhucau}).

\begin{table}[htbp]
    \centering
    \caption{Phân loại dạng nhu cầu của $30.490$ chuỗi cửa hàng--mặt hàng trong
    bộ dữ liệu M5 theo tiêu chí của \textcite{syntetos2005categorization}}
    \label{tab:phanloainhucau}
    \begin{tabular}{lcc}
        \hline
        \textbf{Nhóm} & \textbf{Đặc điểm} & \textbf{Tỷ lệ} \\
        \hline
        Smooth       & $p < 1{,}32$; $CV^2 < 0{,}49$   & $3{,}2\%$  \\
        Erratic      & $p < 1{,}32$; $CV^2 \ge 0{,}49$ & $1{,}6\%$  \\
        Intermittent & $p \ge 1{,}32$; $CV^2 < 0{,}49$ & $75{,}7\%$ \\
        Lumpy        & $p \ge 1{,}32$; $CV^2 \ge 0{,}49$ & $19{,}5\%$ \\
        \hline
    \end{tabular}
\end{table}

Đặc điểm gián đoạn này khiến các phương pháp dự báo và các công thức tồn kho
xây dựng trên giả định nhu cầu liên tục và phân phối chuẩn trở nên kém phù hợp
\parencite{croston1972, syntetos2005accuracy}. Quản lý tồn kho trên một mạng
lưới như vậy đòi hỏi những công cụ ra quyết định có khả năng thích ứng, thay vì
các bộ tham số tĩnh được thiết lập một lần và áp dụng cố định.

\subsection{Thách thức trong việc tối ưu hóa chính sách đặt hàng lại}
\label{subsec:thachthuc}

Việc xây dựng chính sách đặt hàng lại trong mạng lưới đa kho -- đa mặt hàng đối
mặt với bốn nhóm thách thức chính.

\textbf{Thứ nhất, tính bất định kép của nhu cầu và thời gian giao hàng.} Nhu
cầu tiêu thụ biến động ngẫu nhiên theo thời gian, đồng thời thời gian từ lúc
phát lệnh đặt hàng đến lúc hàng thực nhập kho cũng không cố định
\parencite{zipkin2000}. Khi thời gian giao hàng là biến ngẫu nhiên, các đơn
hàng phát ra ở những thời điểm khác nhau có thể về kho không đúng thứ tự, khiến
các công thức tính tồn kho an toàn cổ điển chỉ còn giữ vai trò xấp xỉ
\parencite{silver2016}.

\textbf{Thứ hai, bài toán đánh đổi giữa nhiều nhóm chi phí xung đột.} Mục tiêu
vận hành không quy về một đại lượng duy nhất mà là tổng hợp của nhiều thành
phần có chiều tác động ngược nhau, được tóm tắt trong
Bảng~\ref{tab:chiphi}. Giảm tồn kho làm giảm chi phí lưu kho nhưng làm tăng
rủi ro thiếu hàng; đặt hàng theo lô lớn làm giảm chi phí đặt hàng cố định nhưng
làm tăng chi phí lưu kho và nguy cơ vượt sức chứa \parencite{zipkin2000}.

\begin{table}[htbp]
    \centering
    \caption{Các nhóm chi phí trong bài toán tồn kho đang xét}
    \label{tab:chiphi}
    \begin{tabular}{llc}
        \hline
        \textbf{Thành phần chi phí} & \textbf{Ký hiệu} & \textbf{Chiều tác động} \\
        \hline
        Chi phí lưu kho          & $c_{lk}$    & Tăng khi mức tồn kho cao \\
        Chi phí thiếu hàng       & $c_{th}$    & Tăng khi xảy ra hết hàng \\
        Chi phí đặt hàng cố định & $c_{dh}$    & Tăng theo tần suất đặt hàng \\
        Phạt tràn kho            & $p_{tk}$    & Tăng khi vượt sức chứa kho \\
        Phạt vi phạm mức phục vụ & $\phi_{dv}$ & Tăng khi tỷ lệ đáp ứng dưới ngưỡng \\
        \hline
    \end{tabular}
\end{table}

\textbf{Thứ ba, ràng buộc sức chứa tạo ra sự phụ thuộc giữa các mặt hàng.} Các
mặt hàng cùng nằm trong một kho chia sẻ một dung lượng hữu hạn. Quyết định đặt
hàng của một mặt hàng do đó ảnh hưởng trực tiếp tới chi phí mà các mặt hàng
khác trong cùng kho phải gánh chịu. Bài toán vì vậy không thể tách rời thành
các bài toán độc lập, mà mang bản chất của một bài toán ra quyết định hợp tác
\parencite{oliehoek2016}.

\textbf{Thứ tư, quy mô không gian hành động.} Với $N$ kho và $M$ mặt hàng, mỗi
cặp có $K$ mức đặt hàng khả dĩ, không gian hành động chung là tích Descartes
có kích thước $K^{NM}$. Với cấu hình của khóa luận ($N = 10$, $M = 50$,
$K = 6$), con số này là $6^{500}$ --- vượt xa khả năng xử lý của bất kỳ phương
pháp tìm kiếm tập trung nào, và cũng vượt quá khả năng biểu diễn của các thuật
toán học tăng cường dựa trên hàm giá trị vốn phải liệt kê toàn bộ hành động khả
dĩ \parencite{sutton2018}. Trên phương diện lý thuyết độ phức tạp, bài toán
điều khiển phi tập trung dạng này đã được chứng minh thuộc lớp NEXP-đầy đủ
\parencite{bernstein2002}, nên các lời giải chính xác không khả thi ở quy mô
thực tế.

\subsection{Hạn chế của phương pháp truyền thống và sự cần thiết của học tăng
cường đa tác tử}
\label{subsec:khoangtrong}

Các chính sách tồn kho cổ điển đều được xây dựng trên những tập giả định cụ
thể. Mô hình lượng đặt hàng kinh tế của \textcite{harris1913} giả định nhu cầu
xác định và đều, không có tình trạng thiếu hàng. Chính sách $(s, S)$ được
\textcite{scarf1960} chứng minh tối ưu cho hệ đơn kho với nhu cầu độc lập giữa
các kỳ, chi phí đặt hàng cố định và cơ chế giữ đơn hàng chờ. Mô hình người bán
báo của \textcite{arrow1951} chỉ tối ưu cho một chu kỳ đơn lẻ, không xét tới
lượng hàng đang trên đường vận chuyển. Không mô hình nào trong số đó xử lý được
đồng thời nhu cầu ngẫu nhiên, thời gian giao hàng biến động, ràng buộc sức chứa
dùng chung và yêu cầu duy trì mức phục vụ tối thiểu. Riêng với cơ chế mất đơn
hàng khi thiếu hàng --- vốn phù hợp với bối cảnh bán lẻ --- \textcite{zipkin2008}
chỉ ra rằng bài toán không còn nghiệm giải tích dạng đóng và các chính sách cổ
điển mất tính tối ưu đã được chứng minh.

Học tăng cường cung cấp một hướng tiếp cận khác: thay vì giả định trước dạng
phân phối của nhu cầu và suy ra chính sách bằng giải tích, tác tử học chính
sách trực tiếp thông qua tương tác lặp lại với môi trường
\parencite{sutton2018}. Hướng này đã được kiểm chứng trên các bài toán tồn kho
quy mô nhỏ: \textcite{oroojlooyjadid2022} áp dụng mạng Q sâu cho trò chơi phân
phối bia, còn \textcite{gijsbrechts2022} đánh giá học tăng cường sâu trên ba
lớp bài toán tồn kho gồm mất đơn hàng, hai nguồn cung và đa cấp.

Tuy nhiên, việc áp dụng học tăng cường đơn tác tử cho toàn bộ mạng lưới vấp
phải chính vấn đề quy mô đã nêu ở mục trên. Học tăng cường đa tác tử với cơ chế
chia sẻ tham số cho phép vượt qua trở ngại này \parencite{gupta2017}: mỗi cặp
kho--mặt hàng được xem là một tác tử ra quyết định độc lập trên quan sát cục
bộ theo hướng tiếp cận học độc lập của \textcite{tan1993}, trong khi toàn bộ
các tác tử cùng học chung một bộ trọng số. Nhờ đó, không gian hành động mà mỗi
tác tử phải xử lý thu về kích thước $K$, còn số tham số cần học không tăng theo
quy mô hệ thống. Các nghiên cứu gần đây cho thấy phương pháp học độc lập dựa
trên PPO đạt hiệu năng cạnh tranh trên nhiều môi trường hợp tác
\parencite{dewitt2020ippo, yu2022mappo}.

Khảo sát tài liệu cho thấy phần lớn các nghiên cứu áp dụng học tăng cường cho
quản lý tồn kho hiện nay được tiến hành ở quy mô nhỏ, với một vài kho và một
vài mặt hàng, trên dữ liệu nhu cầu sinh theo phân phối lý thuyết. Số công trình
đồng thời mở rộng tới quy mô hàng trăm cặp kho--mặt hàng bằng cơ chế chia sẻ
tham số, đánh giá tường minh ảnh hưởng của cách phân rã hàm phần thưởng ở quy
mô đó, và hiệu chỉnh môi trường mô phỏng từ dữ liệu bán lẻ thực tế còn hạn chế.
Đây chính là khoảng trống mà khóa luận hướng tới.

\section{Mục tiêu và câu hỏi nghiên cứu}
\label{sec:muctieu}

\subsection{Mục tiêu nghiên cứu}
\label{subsec:muctieu}

Khóa luận đặt ra bốn mục tiêu cụ thể:

\begin{itemize}
    \item Xây dựng môi trường mô phỏng tồn kho đa kho -- đa mặt hàng theo chuẩn
    Gymnasium \parencite{towers2024gymnasium}, mô hình hóa đầy đủ nhu cầu ngẫu
    nhiên, thời gian giao hàng biến động, ràng buộc sức chứa và năm nhóm chi phí
    vận hành; tham số nhu cầu được hiệu chỉnh từ bộ dữ liệu M5 Walmart.

    \item Thiết kế hàm phần thưởng phân rã cục bộ kèm thành phần phạt dựa trên
    tỷ lệ đáp ứng nhu cầu tính trên cửa sổ trượt, bảo đảm tín hiệu học phản ánh
    đúng đóng góp của từng tác tử \parencite{wolpert2001}.

    \item Cài đặt và huấn luyện thuật toán Independent Multi-Agent PPO
    \parencite{schulman2017ppo, dewitt2020ippo} với cơ chế chia sẻ tham số cho
    $500$ cặp kho--mặt hàng, sử dụng ước lượng lợi thế tổng quát
    \parencite{schulman2016gae} và chuẩn hóa lợi thế theo từng cặp.

    \item Đánh giá hiệu năng của mô hình đề xuất trong so sánh với ba chính
    sách tồn kho cổ điển, theo giao thức thực nghiệm đa hạt giống có kiểm định
    thống kê \parencite{henderson2018, agarwal2021}.
\end{itemize}

\subsection{Câu hỏi nghiên cứu}
\label{subsec:cauhoi}

Khóa luận hướng tới trả lời ba câu hỏi:

\begin{enumerate}
    \item \textbf{Câu hỏi 1.} Tác tử IPPO có học được chính sách đặt hàng lại
    cho tổng chi phí vận hành thấp hơn các chính sách tối ưu hóa cổ điển
    (EOQ, $(s, S)$, Newsvendor) trong môi trường đa kho -- đa mặt hàng với nhu
    cầu ngẫu nhiên và thời gian giao hàng biến động hay không, và mức cải thiện
    đó có ý nghĩa thống kê hay không?

    \item \textbf{Câu hỏi 2.} Cơ chế chia sẻ tham số giữa $500$ tác tử có đủ để
    học một chính sách tổng quát hóa được cho các mặt hàng có quy mô nhu cầu
    khác biệt đáng kể hay không?

    \item \textbf{Câu hỏi 3.} Hàm phần thưởng phân rã cục bộ kèm tỷ lệ đáp ứng
    theo cửa sổ trượt có cải thiện quá trình học so với hàm phần thưởng toàn
    cục dùng chung hay không?
\end{enumerate}

\section{Đối tượng và phạm vi nghiên cứu}
\label{sec:phamvi}

\subsection{Đối tượng nghiên cứu}
\label{subsec:doituong}

Đối tượng nghiên cứu của khóa luận là bài toán tối ưu hóa chính sách đặt hàng
lại trong hệ thống tồn kho đa kho -- đa mặt hàng, và thuật toán học tăng cường
đa tác tử Independent Multi-Agent PPO áp dụng cho bài toán này.

\subsection{Phạm vi trong đề tài}
\label{subsec:phamvitrong}

\begin{itemize}
    \item \textbf{Cấu hình hệ thống:} $N = 10$ kho, $M = 50$ mặt hàng, tương
    ứng $500$ cặp ra quyết định; độ dài mỗi tập huấn luyện là $365$ ngày; thời
    gian giao hàng ngẫu nhiên trong khoảng từ một đến ba ngày.

    \item \textbf{Cấu trúc mạng lưới:} các kho hoạt động song song và độc lập
    về nguồn cung, mỗi kho đặt hàng trực tiếp từ nhà cung cấp bên ngoài. Sự
    ràng buộc giữa các mặt hàng phát sinh từ dung lượng kho dùng chung.

    \item \textbf{Nguồn dữ liệu:} bộ dữ liệu bán lẻ M5 Walmart
    \parencite{makridakis2022background}, cụ thể là hai tệp
    \texttt{sales\_train\_evaluation.csv} chứa $30.490$ chuỗi doanh số theo
    ngày trên $1.941$ ngày, và \texttt{calendar.csv} chứa thông tin lịch, ngày
    lễ và các sự kiện trợ cấp thực phẩm SNAP. Khóa luận đồng thời hỗ trợ sinh
    dữ liệu giả lập theo phân phối Poisson phục vụ thử nghiệm đối chứng.

    \item \textbf{Phương pháp đối chứng:} ba chính sách tồn kho cổ điển EOQ
    \parencite{harris1913}, $(s, S)$ \parencite{scarf1960} và Newsvendor
    \parencite{arrow1951}, với tham số được tối ưu bằng tìm kiếm lưới trên
    chính môi trường mô phỏng nhằm bảo đảm tính công bằng của so sánh.

    \item \textbf{Chỉ số đánh giá:} tổng chi phí vận hành, tỷ lệ đáp ứng nhu
    cầu, và cơ cấu các thành phần chi phí lưu kho, thiếu hàng, đặt hàng.
\end{itemize}

\subsection{Giới hạn phạm vi}
\label{subsec:phamvingoai}

Để bảo đảm chiều sâu nghiên cứu trong thời gian thực hiện, những nội dung sau
được xác định nằm ngoài phạm vi khóa luận:

\begin{itemize}
    \item Cấu trúc tồn kho đa cấp (multi-echelon) với quan hệ cung ứng giữa các
    kho theo mô hình của \textcite{clark1960}, cùng các hiện tượng đi kèm như
    hiệu ứng bullwhip \parencite{lee1997}.
    \item Bài toán điều chuyển hàng liên kho và tối ưu hóa tuyến vận chuyển.
    \item Dự báo nhu cầu bằng mô hình học sâu chuỗi thời gian; khóa luận sử
    dụng trực tiếp thông tin nhu cầu lịch sử trong vector quan sát.
    \item Cơ chế truyền thông giữa các tác tử; các tác tử ra quyết định hoàn
    toàn dựa trên quan sát cục bộ.
    \item Triển khai hệ thống thời gian thực hoặc tích hợp với nền tảng quản trị
    doanh nghiệp; sản phẩm của khóa luận là kết quả mô phỏng và phân tích định
    lượng.
\end{itemize}

\section{Phương pháp nghiên cứu}
\label{sec:phuongphap}

Khóa luận kết hợp ba nhóm phương pháp. Về \textbf{lý thuyết}, bài toán được mô
hình hóa dưới dạng quá trình quyết định Markov phi tập trung có quan sát bộ
phận \parencite{bernstein2002, oliehoek2016}, trên nền tảng lý thuyết kiểm soát
tồn kho ngẫu nhiên \parencite{zipkin2000} và học tăng cường
\parencite{puterman1994, sutton2018}. Về \textbf{thực nghiệm}, nhóm xây dựng
môi trường mô phỏng có kiểm thử đơn vị theo chuẩn giao diện Gymnasium
\parencite{brockman2016gym, towers2024gymnasium}, hiệu chỉnh tham số từ dữ liệu
M5, huấn luyện tác tử và thu thập số liệu vận hành. Về \textbf{phân tích}, mọi
so sánh giữa các phương pháp đều được thực hiện trên tối thiểu năm hạt giống
ngẫu nhiên, báo cáo giá trị trung bình kèm độ lệch chuẩn, và được xác nhận bằng
kiểm định thống kê phù hợp trước khi kết luận --- theo khuyến nghị về tính tái
lập trong đánh giá học tăng cường sâu của \textcite{henderson2018} và
\textcite{agarwal2021}.

\section{Ý nghĩa khoa học và thực tiễn của đề tài}
\label{sec:ynghia}

\textbf{Về mặt khoa học}, khóa luận đóng góp một khung mô hình hóa bài toán tồn
kho đa kho -- đa mặt hàng dưới dạng hệ đa tác tử hợp tác quy mô lớn, cùng với
đánh giá thực nghiệm có kiểm định thống kê về ảnh hưởng của cách phân rã hàm
phần thưởng và cơ chế chia sẻ tham số --- hai yếu tố thiết kế thường được sử
dụng nhưng ít được khảo sát tách bạch trong các nghiên cứu trước.

\textbf{Về mặt thực tiễn}, kết quả của khóa luận cung cấp bằng chứng định lượng
về mức độ khả thi của việc thay thế các chính sách đặt hàng theo tham số tĩnh
bằng chính sách học được từ dữ liệu, trên một môi trường mô phỏng được hiệu
chỉnh từ dữ liệu bán lẻ thực tế. Cấu trúc mã nguồn và giao thức đánh giá được
tổ chức theo hướng có thể tái lập, tạo điều kiện cho các nghiên cứu mở rộng
tiếp theo.

\section{Bố cục của khóa luận}
\label{sec:bocuc}

Khóa luận được tổ chức thành năm chương:

\begin{itemize}
    \item \textbf{Chương 1 --- Tổng quan về đề tài:} trình bày bối cảnh, tính
    cấp thiết, mục tiêu, câu hỏi nghiên cứu, phạm vi và phương pháp nghiên cứu.

    \item \textbf{Chương 2 --- Cơ sở lý thuyết và các công trình liên quan:}
    hệ thống hóa lý thuyết kiểm soát tồn kho ngẫu nhiên, các chính sách cổ
    điển, nền tảng học tăng cường và học tăng cường đa tác tử, cùng lược khảo
    các nghiên cứu tiền nhiệm.

    \item \textbf{Chương 3 --- Phân tích và thiết kế mô hình:} chi tiết hóa mô
    hình toán học của môi trường, thiết kế không gian trạng thái, không gian
    hành động và hàm phần thưởng, kiến trúc mạng Actor--Critic chia sẻ tham số
    và quy trình huấn luyện.

    \item \textbf{Chương 4 --- Thực nghiệm và đánh giá kết quả:} trình bày cấu
    hình thực nghiệm, kết quả so sánh với các chính sách cổ điển, các khảo sát
    loại trừ và phân tích độ nhạy.

    \item \textbf{Chương 5 --- Kết luận và hướng phát triển:} tổng kết kết quả
    đạt được, chỉ ra hạn chế và đề xuất hướng nghiên cứu tiếp theo.
\end{itemize}

%
%  GÓI CẦN CÓ TRONG PREAMBLE:
%     \usepackage{amsmath,amssymb,array}
%     \usepackage[hidelinks]{hyperref}
%     \usepackage[style=apa,backend=biber]{biblatex}
%     \addbibresource{tailieuthamkhao.bib}
%
%  LỆNH TRÍCH DẪN: file này dùng biblatex-apa (\parencite, \textcite).
%     - Nếu khóa luận đang dùng natbib/apacite thì thay:
%           \parencite{key}  ->  \citep{key}
%           \textcite{key}   ->  \citet{key}
%
%  SỐ LIỆU M5 trong mục 1.1.1 và 1.3.2 được tính trực tiếp từ hai tệp
%  calendar.csv và sales_train_evaluation.csv (xem Phụ lục nguồn trích dẫn).
% =====================================================================

\chapter{TỔNG QUAN VỀ ĐỀ TÀI}
\label{chap:tongquan}

\section{Bối cảnh nghiên cứu và tính cấp thiết của đề tài}
\label{sec:boicanh}

\subsection{Thực trạng quản lý tồn kho trong mạng lưới bán lẻ đa điểm}
\label{subsec:thuctrang}

Trong hoạt động bán lẻ hiện đại, hàng hóa không được lưu trữ tập trung tại một
điểm duy nhất mà được phân tán trên một mạng lưới gồm nhiều kho hoặc nhiều cửa
hàng đặt tại các khu vực địa lý khác nhau. Mỗi điểm lưu trữ phục vụ một tập
khách hàng riêng, đối mặt với một chuỗi nhu cầu riêng và phải tự ra quyết định
bổ sung hàng hóa cho hàng trăm đến hàng nghìn mặt hàng khác nhau
\parencite{silver2016}. Bài toán vì vậy không còn là việc xác định một mức đặt
hàng đơn lẻ, mà là việc điều phối đồng thời một số lượng rất lớn các quyết định
có ràng buộc lẫn nhau.

Đặc điểm này thể hiện rõ trong dữ liệu bán lẻ thực tế. Bộ dữ liệu M5, công bố
trong cuộc thi dự báo M5 do \textcite{makridakis2022background} tổ chức, ghi
nhận doanh số bán hàng theo ngày của $3.049$ mặt hàng tại $10$ cửa hàng của
Walmart thuộc ba bang California, Texas và Wisconsin, trải trên $1.941$ ngày từ
ngày 29 tháng 01 năm 2011 đến ngày 22 tháng 05 năm
2016\footnote{Số liệu được tính trực tiếp từ hai tệp \texttt{calendar.csv} và
\texttt{sales\_train\_evaluation.csv} sử dụng trong khóa luận. Việc sử dụng bộ
dữ liệu tuân thủ điều khoản sử dụng dữ liệu của cuộc thi cho mục đích nghiên
cứu học thuật phi thương mại \parencite{m5dataset}.}. Tổ hợp cửa hàng và mặt
hàng tạo thành $30.490$ chuỗi thời gian nhu cầu ở cấp độ chi tiết nhất.

Thống kê trên toàn bộ $30.490$ chuỗi này cho thấy hai đặc trưng có ảnh hưởng
trực tiếp tới thiết kế chính sách đặt hàng. Thứ nhất, nhu cầu ở cấp cửa
hàng--mặt hàng có quy mô rất nhỏ và phân bố lệch mạnh: nhu cầu trung bình là
$1{,}13$ đơn vị mỗi ngày, nhưng trung vị chỉ là $0{,}45$ đơn vị, và $73{,}4\%$
số chuỗi có nhu cầu trung bình dưới một đơn vị mỗi ngày. Thứ hai, nhu cầu mang
tính gián đoạn rõ rệt, với tỷ lệ ngày không phát sinh giao dịch trung bình đạt
$68\%$. Áp dụng tiêu chí phân loại của \textcite{syntetos2005categorization} với
các ngưỡng khoảng cách giữa hai lần phát sinh nhu cầu $p = 1{,}32$ và bình
phương hệ số biến thiên $CV^2 = 0{,}49$, có tới $95{,}2\%$ số chuỗi thuộc hai
nhóm \textit{intermittent} và \textit{lumpy}
(Bảng~\ref{tab:phanloainhucau}).

\begin{table}[htbp]
    \centering
    \caption{Phân loại dạng nhu cầu của $30.490$ chuỗi cửa hàng--mặt hàng trong
    bộ dữ liệu M5 theo tiêu chí của \textcite{syntetos2005categorization}}
    \label{tab:phanloainhucau}
    \begin{tabular}{lcc}
        \hline
        \textbf{Nhóm} & \textbf{Đặc điểm} & \textbf{Tỷ lệ} \\
        \hline
        Smooth       & $p < 1{,}32$; $CV^2 < 0{,}49$   & $3{,}2\%$  \\
        Erratic      & $p < 1{,}32$; $CV^2 \ge 0{,}49$ & $1{,}6\%$  \\
        Intermittent & $p \ge 1{,}32$; $CV^2 < 0{,}49$ & $75{,}7\%$ \\
        Lumpy        & $p \ge 1{,}32$; $CV^2 \ge 0{,}49$ & $19{,}5\%$ \\
        \hline
    \end{tabular}
\end{table}

Đặc điểm gián đoạn này khiến các phương pháp dự báo và các công thức tồn kho
xây dựng trên giả định nhu cầu liên tục và phân phối chuẩn trở nên kém phù hợp
\parencite{croston1972, syntetos2005accuracy}. Quản lý tồn kho trên một mạng
lưới như vậy đòi hỏi những công cụ ra quyết định có khả năng thích ứng, thay vì
các bộ tham số tĩnh được thiết lập một lần và áp dụng cố định.

\subsection{Thách thức trong việc tối ưu hóa chính sách đặt hàng lại}
\label{subsec:thachthuc}

Việc xây dựng chính sách đặt hàng lại trong mạng lưới đa kho -- đa mặt hàng đối
mặt với bốn nhóm thách thức chính.

\textbf{Thứ nhất, tính bất định kép của nhu cầu và thời gian giao hàng.} Nhu
cầu tiêu thụ biến động ngẫu nhiên theo thời gian, đồng thời thời gian từ lúc
phát lệnh đặt hàng đến lúc hàng thực nhập kho cũng không cố định
\parencite{zipkin2000}. Khi thời gian giao hàng là biến ngẫu nhiên, các đơn
hàng phát ra ở những thời điểm khác nhau có thể về kho không đúng thứ tự, khiến
các công thức tính tồn kho an toàn cổ điển chỉ còn giữ vai trò xấp xỉ
\parencite{silver2016}.

\textbf{Thứ hai, bài toán đánh đổi giữa nhiều nhóm chi phí xung đột.} Mục tiêu
vận hành không quy về một đại lượng duy nhất mà là tổng hợp của nhiều thành
phần có chiều tác động ngược nhau, được tóm tắt trong
Bảng~\ref{tab:chiphi}. Giảm tồn kho làm giảm chi phí lưu kho nhưng làm tăng
rủi ro thiếu hàng; đặt hàng theo lô lớn làm giảm chi phí đặt hàng cố định nhưng
làm tăng chi phí lưu kho và nguy cơ vượt sức chứa \parencite{zipkin2000}.

\begin{table}[htbp]
    \centering
    \caption{Các nhóm chi phí trong bài toán tồn kho đang xét}
    \label{tab:chiphi}
    \begin{tabular}{llc}
        \hline
        \textbf{Thành phần chi phí} & \textbf{Ký hiệu} & \textbf{Chiều tác động} \\
        \hline
        Chi phí lưu kho          & $c_{lk}$    & Tăng khi mức tồn kho cao \\
        Chi phí thiếu hàng       & $c_{th}$    & Tăng khi xảy ra hết hàng \\
        Chi phí đặt hàng cố định & $c_{dh}$    & Tăng theo tần suất đặt hàng \\
        Phạt tràn kho            & $p_{tk}$    & Tăng khi vượt sức chứa kho \\
        Phạt vi phạm mức phục vụ & $\phi_{dv}$ & Tăng khi tỷ lệ đáp ứng dưới ngưỡng \\
        \hline
    \end{tabular}
\end{table}

\textbf{Thứ ba, ràng buộc sức chứa tạo ra sự phụ thuộc giữa các mặt hàng.} Các
mặt hàng cùng nằm trong một kho chia sẻ một dung lượng hữu hạn. Quyết định đặt
hàng của một mặt hàng do đó ảnh hưởng trực tiếp tới chi phí mà các mặt hàng
khác trong cùng kho phải gánh chịu. Bài toán vì vậy không thể tách rời thành
các bài toán độc lập, mà mang bản chất của một bài toán ra quyết định hợp tác
\parencite{oliehoek2016}.

\textbf{Thứ tư, quy mô không gian hành động.} Với $N$ kho và $M$ mặt hàng, mỗi
cặp có $K$ mức đặt hàng khả dĩ, không gian hành động chung là tích Descartes
có kích thước $K^{NM}$. Với cấu hình của khóa luận ($N = 10$, $M = 50$,
$K = 6$), con số này là $6^{500}$ --- vượt xa khả năng xử lý của bất kỳ phương
pháp tìm kiếm tập trung nào, và cũng vượt quá khả năng biểu diễn của các thuật
toán học tăng cường dựa trên hàm giá trị vốn phải liệt kê toàn bộ hành động khả
dĩ \parencite{sutton2018}. Trên phương diện lý thuyết độ phức tạp, bài toán
điều khiển phi tập trung dạng này đã được chứng minh thuộc lớp NEXP-đầy đủ
\parencite{bernstein2002}, nên các lời giải chính xác không khả thi ở quy mô
thực tế.

\subsection{Hạn chế của phương pháp truyền thống và sự cần thiết của học tăng
cường đa tác tử}
\label{subsec:khoangtrong}

Các chính sách tồn kho cổ điển đều được xây dựng trên những tập giả định cụ
thể. Mô hình lượng đặt hàng kinh tế của \textcite{harris1913} giả định nhu cầu
xác định và đều, không có tình trạng thiếu hàng. Chính sách $(s, S)$ được
\textcite{scarf1960} chứng minh tối ưu cho hệ đơn kho với nhu cầu độc lập giữa
các kỳ, chi phí đặt hàng cố định và cơ chế giữ đơn hàng chờ. Mô hình người bán
báo của \textcite{arrow1951} chỉ tối ưu cho một chu kỳ đơn lẻ, không xét tới
lượng hàng đang trên đường vận chuyển. Không mô hình nào trong số đó xử lý được
đồng thời nhu cầu ngẫu nhiên, thời gian giao hàng biến động, ràng buộc sức chứa
dùng chung và yêu cầu duy trì mức phục vụ tối thiểu. Riêng với cơ chế mất đơn
hàng khi thiếu hàng --- vốn phù hợp với bối cảnh bán lẻ --- \textcite{zipkin2008}
chỉ ra rằng bài toán không còn nghiệm giải tích dạng đóng và các chính sách cổ
điển mất tính tối ưu đã được chứng minh.

Học tăng cường cung cấp một hướng tiếp cận khác: thay vì giả định trước dạng
phân phối của nhu cầu và suy ra chính sách bằng giải tích, tác tử học chính
sách trực tiếp thông qua tương tác lặp lại với môi trường
\parencite{sutton2018}. Hướng này đã được kiểm chứng trên các bài toán tồn kho
quy mô nhỏ: \textcite{oroojlooyjadid2022} áp dụng mạng Q sâu cho trò chơi phân
phối bia, còn \textcite{gijsbrechts2022} đánh giá học tăng cường sâu trên ba
lớp bài toán tồn kho gồm mất đơn hàng, hai nguồn cung và đa cấp.

Tuy nhiên, việc áp dụng học tăng cường đơn tác tử cho toàn bộ mạng lưới vấp
phải chính vấn đề quy mô đã nêu ở mục trên. Học tăng cường đa tác tử với cơ chế
chia sẻ tham số cho phép vượt qua trở ngại này \parencite{gupta2017}: mỗi cặp
kho--mặt hàng được xem là một tác tử ra quyết định độc lập trên quan sát cục
bộ theo hướng tiếp cận học độc lập của \textcite{tan1993}, trong khi toàn bộ
các tác tử cùng học chung một bộ trọng số. Nhờ đó, không gian hành động mà mỗi
tác tử phải xử lý thu về kích thước $K$, còn số tham số cần học không tăng theo
quy mô hệ thống. Các nghiên cứu gần đây cho thấy phương pháp học độc lập dựa
trên PPO đạt hiệu năng cạnh tranh trên nhiều môi trường hợp tác
\parencite{dewitt2020ippo, yu2022mappo}.

Khảo sát tài liệu cho thấy phần lớn các nghiên cứu áp dụng học tăng cường cho
quản lý tồn kho hiện nay được tiến hành ở quy mô nhỏ, với một vài kho và một
vài mặt hàng, trên dữ liệu nhu cầu sinh theo phân phối lý thuyết. Số công trình
đồng thời mở rộng tới quy mô hàng trăm cặp kho--mặt hàng bằng cơ chế chia sẻ
tham số, đánh giá tường minh ảnh hưởng của cách phân rã hàm phần thưởng ở quy
mô đó, và hiệu chỉnh môi trường mô phỏng từ dữ liệu bán lẻ thực tế còn hạn chế.
Đây chính là khoảng trống mà khóa luận hướng tới.

\section{Mục tiêu và câu hỏi nghiên cứu}
\label{sec:muctieu}

\subsection{Mục tiêu nghiên cứu}
\label{subsec:muctieu}

Khóa luận đặt ra bốn mục tiêu cụ thể:

\begin{itemize}
    \item Xây dựng môi trường mô phỏng tồn kho đa kho -- đa mặt hàng theo chuẩn
    Gymnasium \parencite{towers2024gymnasium}, mô hình hóa đầy đủ nhu cầu ngẫu
    nhiên, thời gian giao hàng biến động, ràng buộc sức chứa và năm nhóm chi phí
    vận hành; tham số nhu cầu được hiệu chỉnh từ bộ dữ liệu M5 Walmart.

    \item Thiết kế hàm phần thưởng phân rã cục bộ kèm thành phần phạt dựa trên
    tỷ lệ đáp ứng nhu cầu tính trên cửa sổ trượt, bảo đảm tín hiệu học phản ánh
    đúng đóng góp của từng tác tử \parencite{wolpert2001}.

    \item Cài đặt và huấn luyện thuật toán Independent Multi-Agent PPO
    \parencite{schulman2017ppo, dewitt2020ippo} với cơ chế chia sẻ tham số cho
    $500$ cặp kho--mặt hàng, sử dụng ước lượng lợi thế tổng quát
    \parencite{schulman2016gae} và chuẩn hóa lợi thế theo từng cặp.

    \item Đánh giá hiệu năng của mô hình đề xuất trong so sánh với ba chính
    sách tồn kho cổ điển, theo giao thức thực nghiệm đa hạt giống có kiểm định
    thống kê \parencite{henderson2018, agarwal2021}.
\end{itemize}

\subsection{Câu hỏi nghiên cứu}
\label{subsec:cauhoi}

Khóa luận hướng tới trả lời ba câu hỏi:

\begin{enumerate}
    \item \textbf{Câu hỏi 1.} Tác tử IPPO có học được chính sách đặt hàng lại
    cho tổng chi phí vận hành thấp hơn các chính sách tối ưu hóa cổ điển
    (EOQ, $(s, S)$, Newsvendor) trong môi trường đa kho -- đa mặt hàng với nhu
    cầu ngẫu nhiên và thời gian giao hàng biến động hay không, và mức cải thiện
    đó có ý nghĩa thống kê hay không?

    \item \textbf{Câu hỏi 2.} Cơ chế chia sẻ tham số giữa $500$ tác tử có đủ để
    học một chính sách tổng quát hóa được cho các mặt hàng có quy mô nhu cầu
    khác biệt đáng kể hay không?

    \item \textbf{Câu hỏi 3.} Hàm phần thưởng phân rã cục bộ kèm tỷ lệ đáp ứng
    theo cửa sổ trượt có cải thiện quá trình học so với hàm phần thưởng toàn
    cục dùng chung hay không?
\end{enumerate}

\section{Đối tượng và phạm vi nghiên cứu}
\label{sec:phamvi}

\subsection{Đối tượng nghiên cứu}
\label{subsec:doituong}

Đối tượng nghiên cứu của khóa luận là bài toán tối ưu hóa chính sách đặt hàng
lại trong hệ thống tồn kho đa kho -- đa mặt hàng, và thuật toán học tăng cường
đa tác tử Independent Multi-Agent PPO áp dụng cho bài toán này.

\subsection{Phạm vi trong đề tài}
\label{subsec:phamvitrong}

\begin{itemize}
    \item \textbf{Cấu hình hệ thống:} $N = 10$ kho, $M = 50$ mặt hàng, tương
    ứng $500$ cặp ra quyết định; độ dài mỗi tập huấn luyện là $365$ ngày; thời
    gian giao hàng ngẫu nhiên trong khoảng từ một đến ba ngày.

    \item \textbf{Cấu trúc mạng lưới:} các kho hoạt động song song và độc lập
    về nguồn cung, mỗi kho đặt hàng trực tiếp từ nhà cung cấp bên ngoài. Sự
    ràng buộc giữa các mặt hàng phát sinh từ dung lượng kho dùng chung.

    \item \textbf{Nguồn dữ liệu:} bộ dữ liệu bán lẻ M5 Walmart
    \parencite{makridakis2022background}, cụ thể là hai tệp
    \texttt{sales\_train\_evaluation.csv} chứa $30.490$ chuỗi doanh số theo
    ngày trên $1.941$ ngày, và \texttt{calendar.csv} chứa thông tin lịch, ngày
    lễ và các sự kiện trợ cấp thực phẩm SNAP. Khóa luận đồng thời hỗ trợ sinh
    dữ liệu giả lập theo phân phối Poisson phục vụ thử nghiệm đối chứng.

    \item \textbf{Phương pháp đối chứng:} ba chính sách tồn kho cổ điển EOQ
    \parencite{harris1913}, $(s, S)$ \parencite{scarf1960} và Newsvendor
    \parencite{arrow1951}, với tham số được tối ưu bằng tìm kiếm lưới trên
    chính môi trường mô phỏng nhằm bảo đảm tính công bằng của so sánh.

    \item \textbf{Chỉ số đánh giá:} tổng chi phí vận hành, tỷ lệ đáp ứng nhu
    cầu, và cơ cấu các thành phần chi phí lưu kho, thiếu hàng, đặt hàng.
\end{itemize}

\subsection{Giới hạn phạm vi}
\label{subsec:phamvingoai}

Để bảo đảm chiều sâu nghiên cứu trong thời gian thực hiện, những nội dung sau
được xác định nằm ngoài phạm vi khóa luận:

\begin{itemize}
    \item Cấu trúc tồn kho đa cấp (multi-echelon) với quan hệ cung ứng giữa các
    kho theo mô hình của \textcite{clark1960}, cùng các hiện tượng đi kèm như
    hiệu ứng bullwhip \parencite{lee1997}.
    \item Bài toán điều chuyển hàng liên kho và tối ưu hóa tuyến vận chuyển.
    \item Dự báo nhu cầu bằng mô hình học sâu chuỗi thời gian; khóa luận sử
    dụng trực tiếp thông tin nhu cầu lịch sử trong vector quan sát.
    \item Cơ chế truyền thông giữa các tác tử; các tác tử ra quyết định hoàn
    toàn dựa trên quan sát cục bộ.
    \item Triển khai hệ thống thời gian thực hoặc tích hợp với nền tảng quản trị
    doanh nghiệp; sản phẩm của khóa luận là kết quả mô phỏng và phân tích định
    lượng.
\end{itemize}

\section{Phương pháp nghiên cứu}
\label{sec:phuongphap}

Khóa luận kết hợp ba nhóm phương pháp. Về \textbf{lý thuyết}, bài toán được mô
hình hóa dưới dạng quá trình quyết định Markov phi tập trung có quan sát bộ
phận \parencite{bernstein2002, oliehoek2016}, trên nền tảng lý thuyết kiểm soát
tồn kho ngẫu nhiên \parencite{zipkin2000} và học tăng cường
\parencite{puterman1994, sutton2018}. Về \textbf{thực nghiệm}, nhóm xây dựng
môi trường mô phỏng có kiểm thử đơn vị theo chuẩn giao diện Gymnasium
\parencite{brockman2016gym, towers2024gymnasium}, hiệu chỉnh tham số từ dữ liệu
M5, huấn luyện tác tử và thu thập số liệu vận hành. Về \textbf{phân tích}, mọi
so sánh giữa các phương pháp đều được thực hiện trên tối thiểu năm hạt giống
ngẫu nhiên, báo cáo giá trị trung bình kèm độ lệch chuẩn, và được xác nhận bằng
kiểm định thống kê phù hợp trước khi kết luận --- theo khuyến nghị về tính tái
lập trong đánh giá học tăng cường sâu của \textcite{henderson2018} và
\textcite{agarwal2021}.

\section{Ý nghĩa khoa học và thực tiễn của đề tài}
\label{sec:ynghia}

\textbf{Về mặt khoa học}, khóa luận đóng góp một khung mô hình hóa bài toán tồn
kho đa kho -- đa mặt hàng dưới dạng hệ đa tác tử hợp tác quy mô lớn, cùng với
đánh giá thực nghiệm có kiểm định thống kê về ảnh hưởng của cách phân rã hàm
phần thưởng và cơ chế chia sẻ tham số --- hai yếu tố thiết kế thường được sử
dụng nhưng ít được khảo sát tách bạch trong các nghiên cứu trước.

\textbf{Về mặt thực tiễn}, kết quả của khóa luận cung cấp bằng chứng định lượng
về mức độ khả thi của việc thay thế các chính sách đặt hàng theo tham số tĩnh
bằng chính sách học được từ dữ liệu, trên một môi trường mô phỏng được hiệu
chỉnh từ dữ liệu bán lẻ thực tế. Cấu trúc mã nguồn và giao thức đánh giá được
tổ chức theo hướng có thể tái lập, tạo điều kiện cho các nghiên cứu mở rộng
tiếp theo.

\section{Bố cục của khóa luận}
\label{sec:bocuc}

Khóa luận được tổ chức thành năm chương:

\begin{itemize}
    \item \textbf{Chương 1 --- Tổng quan về đề tài:} trình bày bối cảnh, tính
    cấp thiết, mục tiêu, câu hỏi nghiên cứu, phạm vi và phương pháp nghiên cứu.

    \item \textbf{Chương 2 --- Cơ sở lý thuyết và các công trình liên quan:}
    hệ thống hóa lý thuyết kiểm soát tồn kho ngẫu nhiên, các chính sách cổ
    điển, nền tảng học tăng cường và học tăng cường đa tác tử, cùng lược khảo
    các nghiên cứu tiền nhiệm.

    \item \textbf{Chương 3 --- Phân tích và thiết kế mô hình:} chi tiết hóa mô
    hình toán học của môi trường, thiết kế không gian trạng thái, không gian
    hành động và hàm phần thưởng, kiến trúc mạng Actor--Critic chia sẻ tham số
    và quy trình huấn luyện.

    \item \textbf{Chương 4 --- Thực nghiệm và đánh giá kết quả:} trình bày cấu
    hình thực nghiệm, kết quả so sánh với các chính sách cổ điển, các khảo sát
    loại trừ và phân tích độ nhạy.

    \item \textbf{Chương 5 --- Kết luận và hướng phát triển:} tổng kết kết quả
    đạt được, chỉ ra hạn chế và đề xuất hướng nghiên cứu tiếp theo.
\end{itemize}
